\documentclass{jlreq}
\usepackage{amsmath}
\usepackage{amssymb}
\usepackage{amsthm}
\usepackage{graphicx}
\usepackage{geometry}
\usepackage{caption}
\pagestyle{empty}
\geometry{left=15mm, right=15mm, top=10mm, bottom=10mm}

\begin{document}

\begin{center}
    \LARGE{\textbf{教育実習の振り返り用発表レジュメ}} \footnote{\LaTeX で書いてます．}
\end{center}

\section{実習のプロフィール}
\subsection{所属・名前} 
理工学部　数学科　竹野拓夢

\subsection{実習校・実習期間}
戸田市立笹目中学校　9/1～9/22　$\left( 3週間 \right)$
\subsection{実習教科・担当学年・教壇実習回数}
数学　中学1年生　26回\footnote{中学校なので道徳の授業がありますが，授業変更の関係でできませんでした．} $\left( 1組：7回，2組：7回，3組：6回，4組：6回 \right)$

\section{実習校の特徴}
\subsection{学校の特徴}
一般的な公立中学校であるがICT機器の普及がすさまじく，生徒は一人一台パソコン所有，教室はホワイトボードかつプロジェクター完備，
さらに電子黒板まで設置されているという徹底ぶり．\footnote{戸田市は財政力が豊かであり，全国でも数少ない地方交付税交付金をもらってない自治体である．
だからできることなのかもしれない．}\par
また，部活熱心な生徒が多いように感じた．ただ，教員の働き方改革によるものなのか，毎週月曜日と木曜日は部活がなく，また朝練も廃止されていた．
限られた練習環境ではあったものの，新人戦に向けて一生懸命練習している姿が印象的だった．

\subsection{生徒の特徴}
1年生はとても元気で素直な生徒が多かった．授業を一生懸命受けている姿が印象的だった．
授業を担当した4クラスはどこも生徒仲が良くて，話し合い活動では積極的に活動する様子が見られた．
また，各クラスに外国籍の生徒が少なくとも一人はいた．\footnote{ほとんどは日本語が話せたので特に問題はなかったが，
一人だけ日本語支援が必要な生徒がいたので，その人には英語で対応した．}クラスにはなじめている様子だった．\par
実習期間中に台風が直撃し，下校の時間に雷雨が激しかった影響で生徒を学校に一時待機させる事態が起きた．その際に，1年5組に入って生徒を見守ることになった．
授業を持っていなかったので，初めて会ったがたくさん話をすることができた．
普段関わることのない生徒との関わりができたし，このようなイレギュラー時の教員や学校の対応を経験できたのは良かった．
\footnote{午後6時ごろには雨も雷もやんだので，保護者に迎えに来てもらうという事態は避けられた．}

\subsection{担当ホームルームの特徴}
担当ホームルームは1年1組．男女の仲がとても良くて，給食の時間や休み時間には隔てなく話す様子が見られた．毎日「やりとり帳」という名の一言日記を生徒が書いたものをチェックしてコメントを
していた．そこには，同じクラスの生徒同士で越谷レイクタウン\footnote{埼玉県越谷市にある日本一大きなショッピングモール．埼玉県の中高生にとっては定番のお出かけスポット．}
に遊びに行ったことや，塾があって疲れた，いつも通りの日だった，などいろんなことが書かれていてとても面白かった．
また，私に対して何か聞きたいことがあれば自由に書いてくださいと言ったところ，たくさんの質問が来た．\footnote{特に印象に残っているのは，一番好きな素数は何ですか？，結婚願望はありますか？，
好きな人ができたらどうしたらいいですか？，です．}なかなか生徒全員と同じ密度でコミュニケーションを取るのは難しかったので，こういうツールを活用して生徒理解を深めることができてよかった．

\newpage

\section{実習を終えて見いだされた自己の課題について}

\subsection*{A. 教科指導について　「基礎学力の徹底」}

\subsection{直面した困難}
私が担当した範囲が，方程式の導入から方程式の基本的な計算を扱うところまでだった．方程式が解けるようになるには，「等式の性質」
\footnote{等式の性質とは，簡単に言うと等式の両辺に同じ数を足しても引いてもかけても割っても等式が成り立つ，という性質のことである．}
や「移項」\footnote{移項とは，一方の辺にある項を他方の辺に符号を変えて移動させることである．以下は，$+7$と$2x$を
それぞれ右辺と左辺に移項している．
\begin{align*}
    3x+7&=2x-9 \\
    3x-2x&=-9-7 \\
    x&=-16
\end{align*}}といった
基本的な事柄を理解することも大事であるが，それ以前の単元である「正負の数」と「文字式」で学んだことの定着が必要不可欠である．
机間指導をしていると，今までに学習したこれらの単元の内容が定着していない生徒が多くみられた．特に多かった間違いが3つあり，1つ目は負の数に関する計算ミス，2つ目は
同類項の計算ミス，3つ目は分配法則での計算ミスであった．さらに，小数・分数の計算などといった算数で学習する計算に苦戦する生徒もいた．
以下の表\ref{tab:mistakes}参照．\footnote{数字は変えてます．}
\vspace{3mm}
\begin{center}
    \captionsetup{font=bf}
    \captionof{table}{よく見られた間違いとその原因}
    \label{tab:mistakes}
    \renewcommand{\arraystretch}{1.3}
\begin{tabular}{|c|c|c|}
    \hline
    \textbf{生徒の誤答例} & \textbf{正答} & \textbf{間違いの原因}\\
    \hline
    $(-4)\times(-8) = -32$ & $32$ & 負の数同士の積の符号が違っている．\\
    $-3-9 = -6$ & $-12$ & $9-3=6$ と混同している．\\
    $5x+2 = 7x$ & $5x+2$ & 同類項でない項を足している． \\
    $3(x+4) = 3x+4$ & $3x+12$ & $4$に$3$を分配し忘れている． \\
    $-5(x+7) = -5x+35$ & $-5x-35$ & 符号の処理が適切にできていない．\\
    \hline
\end{tabular}
\end{center}
\vspace{5mm}

\subsection{とった対応}
生徒によって間違いの傾向が異なるので，机間指導を通して個別に対応した．早く終わった生徒に対して，周りに分からない人がいたら
教えてあげるよう指示を出して，生徒同士で学びあえる環境を作った．特に目立った間違いや，重要な事柄については全体に向けて注意した．

\section{最後に}
最近の学校教育は，生徒が主体的に学べる環境を作ることや，グループワークなどの活動を取り入れること，生徒に深い思考を促すような授業をすることが望まれている．
そのこと自体を否定するつもりは全くないし，主体的に学ぶことでより理解が深まるのでとても大切なことであると思う．しかし，それによってあまりにも「基礎的な知識及び技能」が犠牲になっていると感じた．
基本的なことが分かっていないと思考のしようがないし，深い学びを得ることもできない．
数学を本当に理解しよう\footnote{そもそも，教師が数学を本質から理解しているのかどうかは疑問ではある．}と思ったら，基礎的な計算は無意識に手が動くくらいになるまでに演習を重ねないと，この先で必ずつまずくことになる．
そのことを今の学校教育の中で生徒に伝えるにはどうしたらいいのか，考えるきっかけとなった教育実習だった．
\footnote{英語の授業を見学したが，そこでは文法の説明は一切なくてひたすら話すという活動を通して表現を覚えるというものであった．
果たしてそれで英語が話せるようになるのだろうか．}
\end{document}